Papers 4--9 of the VulnPrio research sequence developed components for
hygiene-augmented vulnerability prioritization --- a calibrated scorer,
a deployable online recalibration recipe, a capacity-arrival operating
regime characterisation, and a closed-loop signal-exhaustion theorem.
None of these papers addresses the operational integration question:
\emph{can the components be assembled into an autonomous remediation
framework with pre-registered safety bounds?}

This paper builds and evaluates that framework. \textbf{AutoHeal} is
a six-stage closed-loop pipeline (detect, score, triage, plan, act,
verify+rollback, learn) parameterised by Papers 4--9's findings and
locked under pre-registered safety bounds. We evaluate it on a
2{,}700-row frozen sweep across $K \in \{50, 100, 200\}$, $\lambda = 3$,
12 windows, 25 evaluation seeds, against a human-in-loop baseline and
a fixed-policy control, with CVE attributes drawn from the real
EPSS/KEV public corpus from Paper~4~\S\ref{sec:external_validity}.

The pre-registration locked four hypotheses (H1--H4). \textbf{The
mixed-outcome rejection pattern is the substantive finding.}

\textbf{H1 (Coverage) — supported.} AutoHeal eventually remediates
$97.8\%$ of high-EPSS pairs by Window 6 at both $K = 50$ and
$K = 100$. The core purpose of autonomous remediation is achieved.

\textbf{H2 (Safety) — rejected.} The pre-registered safety bound
($\leq 5\%$ rollback rate per window) is violated at 300 of 900
post-W1 windows. The hard-stop fires 134 times across the sweep;
cascading failures detected three times. \textbf{The safety
mechanism works as designed, but it fires more often than the
pre-registered tolerance contemplated.}

\textbf{H3 (MTTR Reduction) — not analysable.} An instrumentation
issue --- disclosure window tracked at remediation rather than
detection --- produces degenerate MTTR estimates (mean $0.0$ for both
strategies). We report the issue honestly per the pre-registration
stop rule and qualify the abstract accordingly. A corrected
instrumentation is pre-registered for the follow-up evaluation.

\textbf{H4 (Per-Pair Dominance over Human-in-Loop) — rejected.}
AutoHeal beats Human-in-loop in only $2\%$ of (cell, seed, window)
triples. The mechanism is structural: at the pre-registered
conservative triage threshold ($0.80$ score for AUTO classification),
the AUTO bucket averages $\approx 27$ pairs per window at $K = 50$,
which is comparable to Human-in-loop's unconditional 30. The
safety hard-stops further reduce AutoHeal's per-window throughput.
Under the pre-registered configuration, AutoHeal and Human-in-loop
deliver comparable coverage; AutoHeal trades equality for
safety-bounded automation rather than for throughput.

The deployable conclusion is honestly conditional:
\textbf{at the pre-registered conservative thresholds, AutoHeal
achieves automation parity with human-in-loop, not throughput
superiority.} A future evaluation with a pre-registered less
conservative triage threshold (or with a fixed instrumentation
for MTTR) could change H3 and H4. The H1 success and the H2
safety-bound rejection are robust to those changes.

All results are bounded to the synthetic EEHDA evaluation context
with real CVE attribute distributions. External validation on real
fleet telemetry, with the framework's pre-registered safety bounds
preserved, is the most consequential remaining direction.
